% Converted from README.md
\subsection{Primer: Maksimum sekvence}\label{primer-maksimum-sekvence}

Primer \texttt{maximum.dfy} prikazuje metod koji pronalazi najveći element neprazne sekvence. Preduslov \texttt{requires\ values\ !=\ []} isključuje prazan ulaz, jer za praznu sekvencu ne postoji element koji bi mogao da bude maksimum.

\codefile{maximum.dfy}{text}

Postuslovi opisuju rezultat bez oslanjanja na konkretnu implementaciju:

\begin{itemize}
\tightlist
\item
  \texttt{max\ in\ values} zahteva da rezultat bude jedan od elemenata ulazne sekvence;
\item
  kvantifikator \texttt{forall} zahteva da nijedan element sekvence nije veći od rezultata.
\end{itemize}

Petlja održava iste dve ideje kroz invarijante. Prva invarijanta kaže da trenutni kandidat za maksimum i dalje dolazi iz ulazne sekvence. Druga kaže da su svi do sada obrađeni elementi manji ili jednaki kandidatu. Na kraju petlje obrađeni su svi elementi, pa invarijante postaju dovoljne za dokaz postuslova.

Lema \texttt{MaximumIsUnique} formalizuje da dva broja koja zadovoljavaju uslove maksimuma nad istom sekvencom moraju biti jednaka. Takve leme se koriste kada je pomoćno svojstvo jasno matematički, ali ga je korisno imenovati da bi se koristilo u drugim dokazima.

Primer se proverava komandom:

\begin{minted}{bash}
$ dafny verify maximum.dfy
\end{minted}
